\section{Materiais e Métodos}
\label{sec:material_metodos}

A metodologia compreendeu a construção do nariz eletrônico, o delineamento
experimental, a aquisição dos sinais do sistema solo--planta, o
pré-processamento das séries temporais e a classificação das coletas a vácuo.

\subsection{Sistema de nariz eletrônico}
\label{subsec:sistema_nariz}

O nariz eletrônico foi desenvolvido para registrar a assinatura de compostos
voláteis.
Foi composto por microcontrolador ESP32, seis sensores de gases MQ, sensor de
umidade do solo, sensor BMP280 de temperatura e pressão, bomba de vácuo,
válvulas, tubo de PVC acoplado ao vaso, mangueiras, câmara sensorial de vidro
e estrutura externa de MDF. O ESP32 registrou os sinais e controlou a bomba
e as válvulas; uma interface web acessada por Wi-Fi permitiu iniciar e
acompanhar as aquisições.

\begin{figure}[H]
    \centering
    \includegraphics[width=0.30\textwidth]{figures/completo.jpeg}
    \caption{Estrutura externa do nariz eletrônico.}
    \label{fig:estrutura_nariz_eletronico}
\end{figure}

\subsubsection{Matriz de sensores de gases}
\label{subsubsec:matriz_sensores}

A matriz foi constituída por MQ-2, MQ-3, MQ-7, MQ-8, MQ-135 e MQ-138,
posicionados na câmara de vidro sob o mesmo fluxo de amostra. Esses sensores
semicondutores de óxido metálico apresentam sensibilidades parcialmente
sobrepostas. Por isso, a classificação utilizou o padrão conjunto das seis
respostas, sem atribuir a um sensor isolado a identificação química de um
composto. As curvas de sensibilidade dos fabricantes foram empregadas apenas
para descrever as respostas nominais dos modelos MQ.

\begin{figure}[H]
    \centering
    \includegraphics[width=0.30\textwidth]{figures/matriz.png}
    \caption{Matriz dos seis sensores de gases na câmara sensorial.}
    \label{fig:matriz_sensores}
\end{figure}

\subsection{Delineamento experimental e material vegetal}
\label{subsec:delineamento}

O experimento foi conduzido na Estação de Pesquisas Agronema, em Grandes
Rios, Paraná, Brasil, nas coordenadas $24{,}186^{\circ}$ S e
$51{,}399^{\circ}$ O. Foram utilizados dez vasos com plantas de soja,
distribuídos em dois grupos experimentais de cinco vasos cada: controle e
inoculado. Os vasos constituíram as unidades biológicas do experimento.

As aquisições começaram aproximadamente 35 dias após a inoculação. A base
utilizada para a análise com coleta a vácuo contém 36 coletas.




\subsection{Vasos saudáveis e vasos inoculados com \textit{Heterodera glycines}}
\label{subsec:classes_experimentais}

Conforme o delineamento experimental, os vasos 1 a 5 formaram o grupo
controle e os vasos 6 a 10 formaram o grupo inoculado. Os vasos inoculados
receberam aproximadamente 100 mil juvenis de segundo estádio (J2) de
\textit{H. glycines} por vaso.

A Figura \ref{fig:vasos_experimento} apresenta os vasos
utilizados no experimento, divididos entre os grupos
controle e inoculado.

\begin{figure}[H]
    \centering
    \includegraphics[width=0.55\textwidth]{figures/vasos_experimento.png}
    \caption{Vasos de soja utilizados no experimento,
    divididos entre os grupos controle (sem nematoide)
    e inoculado (com \textit{H. glycines}).}
    \label{fig:vasos_experimento}

    \vspace{0.1cm}
    {\footnotesize Fonte: Elaborado pelos autores (2026).}
\end{figure}

Na base de classificação, a \textbf{classe 0} identifica
o grupo inoculado, ou com nematoide, e a \textbf{classe 1}
identifica o grupo controle, ou saudável. Assim, ao longo
do artigo, ``recall da classe com nematoide'' corresponde
ao recall da classe 0.


\subsection{Protocolo de aquisição de dados por coleta a vácuo}
\label{subsec:protocolo_aquisicao}

O tubo de PVC foi acoplado ao vaso para delimitar o espaço de amostragem
(Figura~\ref{fig:vasos_experimentos}). A bomba conduziu o ar do
\textit{headspace} pela tubulação até a câmara dos sensores. Uma execução
completa do ciclo, registrada sob um identificador próprio de vaso e dia, foi
definida como uma coleta. Em cada leitura, o ESP32 registrou o tempo decorrido,
os seis sinais MQ, a leitura do sensor de solo, a temperatura e a pressão.

\begin{figure}[H]
    \centering
    \includegraphics[width=0.55\textwidth]{figures/vasos_experimentos.png}
    \caption{Protótipo posicionado sobre o vaso de soja durante a coleta a vácuo.}
    \label{fig:vasos_experimentos}
    \vspace{0.1cm}
    {\footnotesize Fonte: Elaborado pelos autores (2026).}
\end{figure}

A Figura~\ref{fig:ciclo_coleta} apresenta as etapas de operação do sistema
com e sem controle de pressão. Neste estudo, foram analisadas as coletas
realizadas com controle de pressão.

\begin{figure}[H]
    \centering
    \includegraphics[width=\textwidth]{figures/fluxo_en.png}
    \caption{Fluxograma dos ciclos de coleta com e sem controle de pressão
    e esquema do sistema de amostragem.}
    \label{fig:ciclo_coleta}
    \vspace{0.1cm}
    {\footnotesize Fonte: Elaborado pelos autores (2026).}
\end{figure}

O ciclo de referência, ilustrado na Figura~\ref{fig:ciclo_coleta},
compreendeu aquecimento dos sensores, aproximadamente entre 0 e 180~s; uma
transição operacional até cerca de 200~s; amostragem do \textit{headspace},
aproximadamente entre 200 e 800~s; e desaturação ou recuperação,
aproximadamente entre 800 e 900~s. Esses intervalos descrevem o protocolo de
referência, não limites rígidos aplicados a todas as séries: a duração efetiva
das coletas variou, e parte dos registros terminou antes de 900~s. O controle
da bomba e das válvulas foi realizado pelo ESP32 durante o ciclo.

A condição de solo seco ou úmido foi definida pelo índice operacional
\texttt{Soil\_indice\_0\_1}, obtido pela normalização mínimo--máximo da leitura
do sensor de umidade do solo na base consolidada, conforme a
Equação~\eqref{eq:indice_solo}:
\begin{equation}
    I_{\mathrm{solo}} =
    \frac{\mathrm{Soil}-\min(\mathrm{Soil})}
         {\max(\mathrm{Soil})-\min(\mathrm{Soil})},
    \label{eq:indice_solo}
\end{equation}
em que $\mathrm{Soil}$ é a leitura bruta do sensor de umidade do solo, e
$\min(\mathrm{Soil})$ e $\max(\mathrm{Soil})$ são, respectivamente, os valores
mínimo e máximo dessa leitura na base consolidada.

Leituras com $I_{\mathrm{solo}}\leq 0{,}4$ foram denominadas \emph{solo seco};
leituras com $I_{\mathrm{solo}}>0{,}4$, \emph{solo úmido}. No conjunto de
teste, todas as leituras de cada coleta ficaram do mesmo lado desse limiar,
permitindo classificar também a coleta inteira como seca ou úmida. O índice não
representa porcentagem calibrada de água no solo.

\subsection{Pré-processamento dos sinais e janelamento}
\label{subsec:preprocessamento}

Cada planilha da base original representou uma coleta. As leituras foram
associadas ao identificador da coleta, ao dia, ao vaso e à classe. Na
consolidação, foram removidos 15\% dos registros iniciais e 15\% dos
registros finais de cada série para reduzir a contribuição dos transientes
de início e fim do ciclo. Esse recorte proporcional não equivale a impor
os mesmos limites de tempo a todas as coletas: permaneceu a faixa central
de 70\% de cada série, que pode conter trechos fora da amostragem nominal
de 200 a 800~s. As colunas necessárias foram
convertidas para formato numérico, e registros sem valores válidos nessas
colunas foram descartados.

Para remover transientes de pressão, calculou-se
$\Delta P_t=|P_t-P_{t-1}|$ dentro de cada coleta. Quando
$\Delta P_t\geq 0{,}10$~kPa, foram excluídas a leitura do evento e até
30 leituras anteriores e 30 posteriores. Na versão estrita do filtro,
também foram removidas leituras fora da faixa definida pela mediana global
da pressão $\pm 0{,}50$~kPa, calculada sobre a base consolidada. Após a
filtragem estrita, a base de coletas a vácuo continha 63.492 leituras de
36 coletas. Não foi aplicada média móvel nem outra suavização temporal
aos sinais usados pelos classificadores.

Para a modelagem principal, as entradas foram seis sinais MQ corrigidos
estatisticamente, o índice de solo, a temperatura em $^{\circ}$C e a
pressão em kPa. Para cada MQ, ajustou-se somente com as coletas de
treinamento um \texttt{HuberRegressor} que estimou a componente associada
ao índice de solo, à temperatura e à pressão. Essa componente foi
centralizada pela média das previsões no treinamento e subtraída do
sinal original, conforme a Equação~\eqref{eq:compensacao_mq}:
\begin{equation}
    \label{eq:compensacao_mq}
    MQ_{s,t}^{\mathrm{corr}} = MQ_{s,t} -
    \left(\widehat{MQ}_{s,t}^{\mathrm{amb}}-
    \overline{\widehat{MQ}_{s}^{\mathrm{amb}}}_{\mathrm{treino}}\right).
\end{equation}
Nessa equação, $s$ identifica o sensor MQ, $t$ identifica a leitura, e a
barra representa a média das previsões nas leituras de treinamento. Os
``MQ corrigidos'' são, portanto, sinais com compensação \emph{estatística}
da associação linear estimada com essas três variáveis. Eles não resultam
de calibração física para concentração de gases ou umidade relativa do ar.
As seis entradas MQ, o índice de solo, a temperatura e a pressão totalizaram
nove canais por leitura.

Para a X-TCN, as leituras foram ordenadas por coleta e tempo e segmentadas
em janelas de 128 leituras, com passo de 64 leituras e sobreposição nominal
de 50\%, sem atravessar limites entre coletas. Quando necessário, uma
janela adicional foi alinhada ao fim da coleta para incluir as últimas
leituras; sua sobreposição podia superar 50\%. Cada janela teve formato de
$9\times128$ (canais $\times$ leituras).


\subsection{Modelos de aprendizado de máquina}
\label{subsec:modelos}

Foram comparados um classificador arbóreo denominado \emph{Random Forest}
nas tabelas deste manuscrito e uma rede X-TCN, sigla de \emph{Explainable
Temporal Convolutional Network} (rede convolucional temporal explicável)
\cite{chen2026xtcn}. A implementação local do modelo arbóreo usa
\texttt{ExtraTreesClassifier}, com 900 árvores,
\texttt{max\_features="sqrt"}, mínimo de dez amostras por folha,
\texttt{bootstrap=False}, sem pesos de classe e semente 42.


A X-TCN implementada neste trabalho usou três blocos convolucionais
temporais residuais, com 24, 24 e 32 canais, convoluções de núcleo 3 e
dilatações 1, 2 e 4. A saída dos blocos passou por agregação média global
no tempo, \emph{dropout} de 0,20 e camada linear de duas classes. O
treinamento utilizou 60 épocas, lotes de 32 janelas, otimizador AdamW com
taxa de aprendizado $10^{-3}$ e \emph{weight decay} de $10^{-4}$, e
entropia cruzada com pesos inversamente proporcionais à frequência das
classes nas janelas de treinamento. A semente aleatória foi 42. A rede
produziu probabilidades por janela; para avaliar cada leitura, foi usada
a média das probabilidades das janelas que a cobriam, com limiar 0,5 para
a classe 1. A arquitetura implementada não incluiu o mecanismo de atenção
descrito para a X-TCN de \cite{chen2026xtcn}.

Cada classificador foi ajustado uma vez com todas as coletas de treinamento
sob pressão forçada. Solo seco e solo úmido foram dois recortes do mesmo
conjunto de teste, não dois treinamentos independentes.

\subsection{Estratégia de treinamento e teste}
\label{subsec:treinamento_teste}

A divisão foi feita por \emph{coleta}, dentro de cada classe, com
proporção nominal de 70\% para treinamento e 30\% para teste e semente 42.
Todas as leituras e janelas de uma coleta permaneceram em um único conjunto.
Das 36 coletas, 24 foram destinadas ao treinamento (41.377 leituras: 11
coletas da classe 0 e 13 da classe 1) e 12 ao teste (22.115 leituras:
seis coletas de cada classe).

Essa separação evita que janelas da mesma aquisição apareçam nos dois
conjuntos. Entretanto, alguns vasos aparecem em ambos porque foram medidos
em dias distintos. Assim, o teste avalia novas \emph{coletas}, não vasos
biologicamente independentes, e as métricas não demonstram generalização
para plantas ou vasos nunca vistos pelo modelo.

No teste, o recorte seco incluiu 16.692 leituras de oito coletas: 6.507
leituras da classe 0 e 10.185 da classe 1. O recorte úmido incluiu 5.423
leituras de quatro coletas: 4.152 da classe 0 e 1.271 da classe 1. Essa
distribuição desigual motivou o relato da acurácia balanceada junto da
acurácia convencional.

\subsection{Métricas de desempenho}
\label{subsec:metricas}

Foram calculadas acurácia, acurácia balanceada, F1 macro e matrizes de
confusão no conjunto de teste. A acurácia é a fração de leituras
classificadas corretamente; a acurácia balanceada é a média dos recalls
das classes 0 e 1; e a F1 macro é a média simples das medidas F1 das duas
classes. O recall da classe com nematoide corresponde à sensibilidade
para a classe 0, ou seja, à fração de leituras de vasos inoculados
corretamente identificadas. As métricas foram calculadas separadamente
para os recortes seco e úmido do mesmo teste.


\subsection{Auditoria exploratória das coletas}
\label{subsec:auditoria_coletas_loso}

Além do filtro de pressão descrito acima, examinou-se a estabilidade dos
seis sinais MQ em cada coleta. Para o sensor $s$, a coleta $c$ e a leitura
$t>1$, calcularam-se a variação absoluta entre leituras consecutivas e o
limiar de cada sensor conforme a Equação~\eqref{eq:salto_mq_loso}:
\begin{equation}
    d_{s,c,t}=\left|x_{s,c,t}-x_{s,c,t-1}\right|,
    \qquad
    q_s=Q_{0{,}995}\!\left(\{d_{s,c,t}:c,\ t>1\}\right),
    \label{eq:salto_mq_loso}
\end{equation}
em que $x_{s,c,t}$ é o sinal bruto do MQ e $Q_{0{,}995}$ é o percentil
99,5 das variações reunidas de todas as coletas para aquele sensor.
As diferenças foram calculadas dentro de cada coleta, sem conectar o fim
de uma coleta ao início de outra. Definiu-se uma leitura com salto
multissensor quando pelo menos dois dos seis MQ superaram seus respectivos
limiares. O indicador de salto e a fração de leituras sinalizadas na coleta
foram calculados pela Equação~\eqref{eq:ruido_multissensor_loso}:
\begin{equation}
    J_{c,t}=\mathbf{1}\!\left\{
      \sum_{s\in\mathcal{S}_{\mathrm{MQ}}}
      \mathbf{1}\{d_{s,c,t}\geq q_s\}\geq 2
    \right\},
    \qquad
    R_c=\frac{1}{n_c}\sum_{t=1}^{n_c}J_{c,t},
    \label{eq:ruido_multissensor_loso}
\end{equation}
em que $n_c$ é o número de leituras da coleta. A regra exploratória de
qualidade sinalizou coletas com $R_c\geq 0{,}05$. A avaliação visual das
séries foi usada para contextualizar os saltos. O limiar foi aplicado à
coleta individual, de forma que as demais coletas do mesmo vaso permaneceram
elegíveis. Os percentis e a regra de exclusão foram definidos após a
inspeção da base completa; sua utilização constitui análise de
sensibilidade, e não um protocolo de qualidade estabelecido antes da
aquisição.

Separadamente, as previsões da análise inicial de Extra Trees com MQ2,
MQ3, MQ8, MQ135, MQ138 e as três entradas ambientais foram agrupadas por
coleta. A fração de erros de cada coleta foi calculada pela
Equação~\eqref{eq:erro_coleta_loso}:
\begin{equation}
    E_c=\frac{1}{n_c}\sum_{t=1}^{n_c}
    \mathbf{1}\{\widehat y_{c,t}\neq y_c\},
    \label{eq:erro_coleta_loso}
\end{equation}
em que $y_c$ é a classe experimental da coleta. Foram examinadas as
coletas com $E_c>0{,}5$. Na análise de sensibilidade com exclusões,
retiraram-se as quatro coletas identificadas por esse critério e uma
coleta adicional sinalizada pela auditoria multissensor; uma das quatro
também satisfazia o critério de ruído. A seleção por $E_c$ utilizou os
erros do próprio classificador na mesma base, e não demonstra defeito de
aquisição nas coletas selecionadas. Por isso, todas as comparações após
essas exclusões são exploratórias e condicionadas a essa seleção.

\subsection{Validação complementar com coletas agrupadas por vaso}
\label{subsec:loso_complementar}

Para examinar o efeito das coletas repetidas de um mesmo vaso na avaliação,
primeiro foram retiradas somente as cinco coletas selecionadas na análise
de sensibilidade. Nenhum vaso foi excluído: suas outras coletas
permaneceram na base. Em seguida, aplicou-se a validação
\textit{leave-one-vase-out} (LOSO) às coletas restantes, agrupadas pelo
vaso de origem. Em cada rodada, todas as coletas restantes de um mesmo
vaso foram separadas juntas para teste. O agrupamento utilizou a
combinação entre
classe experimental e código de vaso. Essa chave técnica evitou uma colisão
de identificadores na
planilha derivada, que contém registros da classe inoculada com código
numérico de vaso 1; no delineamento físico, os vasos 1 a 5 são controles
e os vasos 6 a 10 são inoculados. Para cada vaso $v$, a separação entre
as coletas de treinamento e de teste seguiu a
Equação~\eqref{eq:particao_loso}:
\begin{equation}
    \begin{aligned}
        \mathcal{C}_v
        &=\{c\in\mathcal{C}:v(c)=v\},\\
        \mathcal{D}^{\mathrm{teste}}_v
        &=\bigcup_{c\in\mathcal{C}_v}\mathcal{D}_c,\\
        \mathcal{D}^{\mathrm{treino}}_v
        &=\bigcup_{c\in\mathcal{C}\setminus\mathcal{C}_v}\mathcal{D}_c,
    \end{aligned}
    \label{eq:particao_loso}
\end{equation}
em que $\mathcal{C}$ é o conjunto de coletas elegíveis, $v(c)$ identifica
o vaso da coleta $c$ e $\mathcal{D}_c$ reúne suas leituras e classes.
Em cada uma das dez rodadas, treinou-se um modelo com as coletas dos
nove outros vasos e testou-se com todas as coletas elegíveis do vaso
reservado. Assim, coletas do mesmo vaso não participaram simultaneamente
do treinamento e do teste em uma mesma rodada, e cada leitura elegível
recebeu uma previsão fora do grupo usado no ajuste. O procedimento
preservou os dez vasos; apenas as cinco coletas retiradas antes do LOSO
deixaram de contribuir com leituras.

Para Extra Trees, usaram-se sinais MQ brutos, sem a compensação estatística
da Equação~\eqref{eq:compensacao_mq}, e, quando previstos na configuração,
o índice operacional do solo, a temperatura e a pressão. O classificador
foi reajustado em cada rodada com 900 árvores,
\texttt{max\_features="sqrt"}, \texttt{min\_samples\_leaf=10},
\texttt{bootstrap=False} e semente 42. A previsão por leitura foi a classe
com probabilidade estimada de classe 1 igual ou superior a 0,5.

Como comparação complementar, \textit{Random Forest} e a rede temporal
descrita anteriormente foram reajustados em cada rodada LOSO em duas
configurações fixas: MQ2, MQ3, MQ8, MQ135 e MQ138, isoladamente ou
acompanhados das três variáveis ambientais. O Random Forest usou 300
árvores, \texttt{class\_weight="balanced"} e semente 42. Na rede temporal,
a padronização das entradas foi estimada apenas no treinamento de cada
rodada. As janelas mantiveram 128 leituras e passo de 64, com cinco ou oito
canais conforme a configuração; nenhuma janela atravessou o limite de
uma coleta. Para uma leitura $i$ coberta pelo conjunto $\mathcal{W}(i)$
de janelas, a probabilidade da classe 1 e a previsão final foram obtidas
pela Equação~\eqref{eq:agregacao_janelas_loso}:
\begin{equation}
    \widehat p_i(1)=
    \frac{1}{|\mathcal{W}(i)|}
    \sum_{w\in\mathcal{W}(i)}\widehat p_w(1),
    \qquad
    \widehat y_i=\mathbf{1}\{\widehat p_i(1)\geq 0{,}5\}.
    \label{eq:agregacao_janelas_loso}
\end{equation}
Os modelos foram comparados por leitura após reunir as previsões de todas
as rodadas. Essa avaliação não corresponde à partição por coleta de 70\%/30\%
descrita anteriormente. A implementação \texttt{ExtraTreesClassifier}
descrita para a análise anterior é tecnicamente Extra Trees; nesta
comparação, \textit{Random Forest} designa especificamente a implementação
\texttt{RandomForestClassifier}. A rede temporal é identificada aqui como
TCN, pois a implementação usada não possui módulo de atenção ou de
explicabilidade integrado.

\subsection{Busca combinatória das entradas}
\label{subsec:busca_combinatoria_loso}

A busca de entradas foi realizada com Extra Trees e a mesma divisão LOSO.
As nove variáveis candidatas estão definidas na
Equação~\eqref{eq:conjunto_entradas_loso}:
\begin{equation}
    \mathcal{F}=\{\mathrm{MQ2},\mathrm{MQ3},\mathrm{MQ7},\mathrm{MQ8},
    \mathrm{MQ135},\mathrm{MQ138},I_{\mathrm{solo}},T,P\}.
    \label{eq:conjunto_entradas_loso}
\end{equation}
Foram avaliados todos os subconjuntos não vazios $A\subseteq\mathcal{F}$.
Cada uma das nove variáveis pode estar presente ou ausente, produzindo
$2^9$ subconjuntos. Um deles é o conjunto vazio, sem qualquer variável
de entrada para o classificador; por isso, ele foi excluído da busca.
O número de configurações avaliadas foi calculado pela
Equação~\eqref{eq:combinacoes_loso}:
\begin{equation}
    \left|\{A\subseteq\mathcal{F}:A\neq\emptyset\}\right|
    =2^{|\mathcal{F}|}-1=2^9-1=511.
    \label{eq:combinacoes_loso}
\end{equation}
Solo, temperatura e pressão variaram separadamente nessa busca,
inclusive nas configurações sem sensores MQ. O BMP280 fornece
temperatura e pressão como duas variáveis, embora seja um único componente
físico; portanto, a contagem de configurações acima se refere a entradas
do modelo, não a componentes de hardware. Se temperatura e pressão forem
tratadas como um bloco indivisível, passam a existir oito componentes
selecionáveis: seis sensores MQ, o sensor de solo e o BMP280. Cada
componente pode estar presente ou ausente, gerando $2^8$ combinações.
Excluindo a combinação vazia, sem nenhum componente, restam
$2^8-1=255$ configurações de componentes possíveis.

Para cada subconjunto $A$, o mesmo classificador e as mesmas dez partições
por vaso foram usados. A ordenação adotou a acurácia por leitura reunida
entre as rodadas; empates foram resolvidos por menor número total de
variáveis, depois por maior número de MQ e, por fim, por ordem alfabética
dos nomes das entradas. Random Forest e rede temporal não participaram
dessa busca exaustiva. A escolha da configuração e sua avaliação ocorreram
nos mesmos vasos, sem validação aninhada nem conjunto biológico externo;
portanto, a ordenação das configurações é exploratória.

\subsection{Cálculo das métricas na validação por vaso}
\label{subsec:metricas_loso_complementar}

Com as previsões das rodadas LOSO reunidas, a acurácia por leitura foi
calculada conforme a Equação~\eqref{eq:acuracia_loso}:
\begin{equation}
    \mathrm{Acc}=\frac{1}{N}
    \sum_{i=1}^{N}\mathbf{1}\{\widehat y_i=y_i\},
    \label{eq:acuracia_loso}
\end{equation}
em que $N$ é o total de leituras avaliadas. Para cada classe
$k\in\{0,1\}$, o recall, a precisão e a medida F1 foram calculados
pela Equação~\eqref{eq:metricas_classe_loso}:
\begin{equation}
    \mathrm{Recall}_k=\frac{TP_k}{TP_k+FN_k},\qquad
    \mathrm{Prec}_k=\frac{TP_k}{TP_k+FP_k},\qquad
    F1_k=\frac{2\,\mathrm{Prec}_k\,\mathrm{Recall}_k}
    {\mathrm{Prec}_k+\mathrm{Recall}_k},
    \label{eq:metricas_classe_loso}
\end{equation}
em que $TP_k$, $FN_k$ e $FP_k$ são, respectivamente, verdadeiros positivos,
falsos negativos e falsos positivos ao tomar $k$ como classe positiva.
A acurácia balanceada e a média macro de F1 foram calculadas conforme a
Equação~\eqref{eq:metricas_resumo_loso}:
\begin{equation}
    \mathrm{Acc}_{\mathrm{bal}}=
    \frac{\mathrm{Recall}_0+\mathrm{Recall}_1}{2},
    \qquad
    F1_{\mathrm{macro}}=\frac{F1_0+F1_1}{2}.
    \label{eq:metricas_resumo_loso}
\end{equation}
O recall da classe 0 descreve a identificação de leituras do grupo
inoculado; os denominadores das métricas são leituras, não vasos. As
leituras sucessivas de uma mesma coleta e as coletas repetidas do mesmo
vaso não constituem réplicas biológicas independentes. Em particular,
a acurácia reunida atribui mais peso a coletas com maior número de leituras.
